% !TEX program = tectonic
% =============================================================================
%  AULA 2 — Redes sem Fio: Introdução
%  Curso: Redes sem Fio  |  Profa. Anelise Munaretto
%  Convertida de .pdf a Beamer + TikZ; conteúdo atualizado (Wi-Fi 6/6E, 5G).
%  Compilar: tectonic aula2.tex
% =============================================================================
\documentclass[aspectratio=169,11pt]{beamer}
% =============================================================================
%  PREÁMBULO COMÚN para as aulas de Redes sem Fio (Beamer + TikZ)
%  Incluido com % =============================================================================
%  PREÁMBULO COMÚN para as aulas de Redes sem Fio (Beamer + TikZ)
%  Incluido com % =============================================================================
%  PREÁMBULO COMÚN para as aulas de Redes sem Fio (Beamer + TikZ)
%  Incluido com \input{preamble.tex} en cada aula.
%  Paleta de cores e estilos são definidos aqui uma única vez.
% =============================================================================

\usepackage[T1]{fontenc}
\usepackage{amsmath,amssymb,mathtools}
\usepackage{graphicx}
\usepackage{tikz}
\usetikzlibrary{positioning,arrows.meta,calc,backgrounds,decorations.pathmorphing,fit,shadows,shapes.symbols,shapes.geometric}
\usepackage{pgfplots}
\pgfplotsset{compat=1.16}
\usepackage{tabularx,booktabs,multirow,array}
\usepackage[normalem]{ulem}
\usepackage{hyperref}

% ---------- Paleta de cores própria ----------
\definecolor{aneliseAzul}{RGB}{20,60,110}
\definecolor{aneliseVerde}{RGB}{30,140,70}
\definecolor{aneliseLaranja}{RGB}{215,120,20}
\definecolor{aneliseGris}{RGB}{90,90,90}

% ---------- Tema ----------
\usetheme{Madrid}
\usecolortheme{whale}
\setbeamercolor{structure}{fg=aneliseAzul}
\setbeamercolor{title}{fg=aneliseAzul}
\setbeamercolor{frametitle}{fg=aneliseAzul}
\setbeamercolor{block title}{fg=aneliseAzul,bg=aneliseGris!12!white}
\setbeamercolor{block body}{bg=aneliseGris!7!white}
\hypersetup{colorlinks=true,linkcolor=aneliseAzul,urlcolor=aneliseVerde}

% ---------- Comandos auxiliares ----------
\newcommand{\kurose}{Kurose \& Ross, \emph{Computer Networking: a Top-Down Approach}}
\newcommand{\forouzan}{Forouzan, \emph{Data Communications and Networking}}
\newcommand{\stallings}{Stallings, \emph{Wireless Communications \& Networks}}
\newcommand{\tanenbaum}{Tanenbaum, \emph{Computer Networks}}
\newcommand{\rfc}[1]{\href{https://www.rfc-editor.org/rfc/rfc#1.txt}{RFC~#1}}
% -----------------------------------------------------------------------------
%  Estilos TikZ comuns para desenhar redes (posicionamento relativo)
% -----------------------------------------------------------------------------
\tikzset{
  host/.style={draw,rounded corners=2mm,fill=aneliseAzul!15,inner sep=1.5mm,font=\scriptsize},
  rt/.style={draw,rounded corners=1mm,fill=aneliseLaranja!25,inner sep=1mm,font=\scriptsize},
  nube/.style={draw,cloud,aspect=2.2,fill=aneliseGris!15,inner sep=1.5mm,font=\scriptsize},
  el/.style={draw,rounded corners=1.2mm,fill=aneliseGris!18,inner sep=0.8mm,font=\scriptsize},
  enl/.style={thick},
  enlA/.style={thick,-Stealth},
} en cada aula.
%  Paleta de cores e estilos são definidos aqui uma única vez.
% =============================================================================

\usepackage[T1]{fontenc}
\usepackage{amsmath,amssymb,mathtools}
\usepackage{graphicx}
\usepackage{tikz}
\usetikzlibrary{positioning,arrows.meta,calc,backgrounds,decorations.pathmorphing,fit,shadows,shapes.symbols,shapes.geometric}
\usepackage{pgfplots}
\pgfplotsset{compat=1.16}
\usepackage{tabularx,booktabs,multirow,array}
\usepackage[normalem]{ulem}
\usepackage{hyperref}

% ---------- Paleta de cores própria ----------
\definecolor{aneliseAzul}{RGB}{20,60,110}
\definecolor{aneliseVerde}{RGB}{30,140,70}
\definecolor{aneliseLaranja}{RGB}{215,120,20}
\definecolor{aneliseGris}{RGB}{90,90,90}

% ---------- Tema ----------
\usetheme{Madrid}
\usecolortheme{whale}
\setbeamercolor{structure}{fg=aneliseAzul}
\setbeamercolor{title}{fg=aneliseAzul}
\setbeamercolor{frametitle}{fg=aneliseAzul}
\setbeamercolor{block title}{fg=aneliseAzul,bg=aneliseGris!12!white}
\setbeamercolor{block body}{bg=aneliseGris!7!white}
\hypersetup{colorlinks=true,linkcolor=aneliseAzul,urlcolor=aneliseVerde}

% ---------- Comandos auxiliares ----------
\newcommand{\kurose}{Kurose \& Ross, \emph{Computer Networking: a Top-Down Approach}}
\newcommand{\forouzan}{Forouzan, \emph{Data Communications and Networking}}
\newcommand{\stallings}{Stallings, \emph{Wireless Communications \& Networks}}
\newcommand{\tanenbaum}{Tanenbaum, \emph{Computer Networks}}
\newcommand{\rfc}[1]{\href{https://www.rfc-editor.org/rfc/rfc#1.txt}{RFC~#1}}
% -----------------------------------------------------------------------------
%  Estilos TikZ comuns para desenhar redes (posicionamento relativo)
% -----------------------------------------------------------------------------
\tikzset{
  host/.style={draw,rounded corners=2mm,fill=aneliseAzul!15,inner sep=1.5mm,font=\scriptsize},
  rt/.style={draw,rounded corners=1mm,fill=aneliseLaranja!25,inner sep=1mm,font=\scriptsize},
  nube/.style={draw,cloud,aspect=2.2,fill=aneliseGris!15,inner sep=1.5mm,font=\scriptsize},
  el/.style={draw,rounded corners=1.2mm,fill=aneliseGris!18,inner sep=0.8mm,font=\scriptsize},
  enl/.style={thick},
  enlA/.style={thick,-Stealth},
} en cada aula.
%  Paleta de cores e estilos são definidos aqui uma única vez.
% =============================================================================

\usepackage[T1]{fontenc}
\usepackage{amsmath,amssymb,mathtools}
\usepackage{graphicx}
\usepackage{tikz}
\usetikzlibrary{positioning,arrows.meta,calc,backgrounds,decorations.pathmorphing,fit,shadows,shapes.symbols,shapes.geometric}
\usepackage{pgfplots}
\pgfplotsset{compat=1.16}
\usepackage{tabularx,booktabs,multirow,array}
\usepackage[normalem]{ulem}
\usepackage{hyperref}

% ---------- Paleta de cores própria ----------
\definecolor{aneliseAzul}{RGB}{20,60,110}
\definecolor{aneliseVerde}{RGB}{30,140,70}
\definecolor{aneliseLaranja}{RGB}{215,120,20}
\definecolor{aneliseGris}{RGB}{90,90,90}

% ---------- Tema ----------
\usetheme{Madrid}
\usecolortheme{whale}
\setbeamercolor{structure}{fg=aneliseAzul}
\setbeamercolor{title}{fg=aneliseAzul}
\setbeamercolor{frametitle}{fg=aneliseAzul}
\setbeamercolor{block title}{fg=aneliseAzul,bg=aneliseGris!12!white}
\setbeamercolor{block body}{bg=aneliseGris!7!white}
\hypersetup{colorlinks=true,linkcolor=aneliseAzul,urlcolor=aneliseVerde}

% ---------- Comandos auxiliares ----------
\newcommand{\kurose}{Kurose \& Ross, \emph{Computer Networking: a Top-Down Approach}}
\newcommand{\forouzan}{Forouzan, \emph{Data Communications and Networking}}
\newcommand{\stallings}{Stallings, \emph{Wireless Communications \& Networks}}
\newcommand{\tanenbaum}{Tanenbaum, \emph{Computer Networks}}
\newcommand{\rfc}[1]{\href{https://www.rfc-editor.org/rfc/rfc#1.txt}{RFC~#1}}
% -----------------------------------------------------------------------------
%  Estilos TikZ comuns para desenhar redes (posicionamento relativo)
% -----------------------------------------------------------------------------
\tikzset{
  host/.style={draw,rounded corners=2mm,fill=aneliseAzul!15,inner sep=1.5mm,font=\scriptsize},
  rt/.style={draw,rounded corners=1mm,fill=aneliseLaranja!25,inner sep=1mm,font=\scriptsize},
  nube/.style={draw,cloud,aspect=2.2,fill=aneliseGris!15,inner sep=1.5mm,font=\scriptsize},
  el/.style={draw,rounded corners=1.2mm,fill=aneliseGris!18,inner sep=0.8mm,font=\scriptsize},
  enl/.style={thick},
  enlA/.style={thick,-Stealth},
}

\title[Redes sem Fio]{Redes sem Fio\\ Introdução}
\subtitle{Aula 2}
\author[Anelise Munaretto]{Profa. Anelise Munaretto}
\institute[Curso]{Redes sem Fio --- Ciência da Computação}
\date{2026}

\begin{document}

\begin{frame}[plain]
  \titlepage
\end{frame}

% ---------------------------------------------------------------------------
\begin{frame}{Objetivos de aprendizagem}
  Ao final desta aula, você será capaz de:
  \begin{itemize}
    \item \textbf{Explicar} as vantagens e desvantagens das redes sem fio;
    \item \textbf{Analisar} os mecanismos de propagação (reflexão, difração, dispersão) e seus efeitos no enlace;
    \item \textbf{Calcular} o orçamento de enlace (perda de percurso e potência recebida);
    \item \textbf{Comparar} as técnicas de espalhamento espectral (FHSS, DSSS, OFDM);
    \item \textbf{Avaliar} os problemas de acesso ao meio em redes \emph{ad hoc} (terminal escondido e terminal exposto).
  \end{itemize}
  \begin{alertblock}{Pré-requisitos}
    Aula 1 --- Internet: camadas TCP/IP, endereçamento IP, camada de enlace e comutação de pacotes.
  \end{alertblock}
\end{frame}

\section{Vantagens}

\begin{frame}{Vantagens das redes sem fio (1/2)}
  \begin{itemize}
    \item \textbf{Redução no tempo de instalação}
    \item \textbf{Fácil planejamento}
      \begin{itemize}
        \item não exige cabeamento prévio
        \item redes \emph{ad hoc} não exigem planejamento algum
      \end{itemize}
    \item \textbf{Instalação em áreas de difícil cabeamento}
      \begin{itemize}
        \item edifícios históricos, com presença de amianto, etc.
      \end{itemize}
    \item \textbf{Redes temporárias}: feiras, exposições, etc.
    \item \textbf{Maior confiabilidade e robustez}
      \begin{itemize}
        \item pode resistir a desastres
      \end{itemize}
  \end{itemize}
\end{frame}

\begin{frame}{Vantagens das redes sem fio (2/2)}
  \begin{itemize}
    \item \textbf{Permite a ``mobilidade'' dos equipamentos}
    \item \textbf{Estética} (ausência de cabos)
    \item \textbf{Flexibilidade}
      \begin{itemize}
        \item transmissores e receptores podem ser colocados em qualquer lugar (dentro de dispositivos, muros, etc.)
      \end{itemize}
    \item \textbf{Diminuição de custos de infraestrutura}
      \begin{itemize}
        \item cabeamento, conectores, etc.
      \end{itemize}
  \end{itemize}
  \begin{center}
    \begin{tikzpicture}[font=\scriptsize]
      \node[host](a){laptop};
      \node[host,right=1.4cm of a](b){tablet};
      \node[host,right=1.4cm of b](c){smartphone};
      \node[nube,above=1.4cm of b](ap){AP 802.11};
      \draw[thick] (a) -- (ap);
      \draw[thick] (b) -- (ap);
      \draw[thick] (c) -- (ap);
      \node[right=1.6cm of c](lbl){\scriptsize\textbf{mobilidade sem cabos}};
    \end{tikzpicture}
  \end{center}
\end{frame}

\section{Desvantagens}

\begin{frame}{Desvantagens e dificuldades (1/2)}
  \begin{itemize}
    \item \textbf{Baixa Qualidade de Serviço (QoS)}
      \begin{itemize}
        \item banda passante limitada
        \item altas taxas de erro (da ordem de $10^{-4}$)
        \item atrasos (controle de acesso ao meio distribuído)
      \end{itemize}
    \item \textbf{Restrições no uso de frequências}
      \begin{itemize}
        \item regulação governamental (ANATEL, FCC, ETSI)
      \end{itemize}
    \item \textbf{Interferência do sinal de rádio}
      \begin{itemize}
        \item vulnerabilidade a ruídos atmosféricos e transmissões de outros sistemas
        \item propagação por múltiplos caminhos
      \end{itemize}
  \end{itemize}
  \note{Reforce que o meio sem fio tem BER da ordem de $10^{-4}$ (contra $10^{-7}$ do cabeado); antecipe o SINR como medida central da qualidade.}
\end{frame}

\begin{frame}{Desvantagens e dificuldades (2/2)}
  \begin{itemize}
    \item \textbf{Consumo de energia} (equipamentos móveis)
    \item \textbf{Desempenho}
    \item \textbf{Segurança}
      \begin{itemize}
        \item baixa privacidade
        \item intrusão, congestionamento da estação-base
      \end{itemize}
    \item \textbf{Aspectos ligados à instalação}
      \begin{itemize}
        \item cobertura, testes de propagação
      \end{itemize}
    \item \textbf{Riscos à saúde} (debate sobre exposição RF)
  \end{itemize}
  \begin{exampleblock}{Nota atualizada}
    WPA3 (2018) e WPA3-Enterprise resolveram boa parte dos problemas clássicos de segurança; a privacidade (captura de identidade) segue sendo um desafio.
  \end{exampleblock}
\end{frame}

\section{Introdução às comunicações sem fio}

\begin{frame}{Comunicação sem fio: perspectiva histórica}
  \begin{itemize}
    \item Usada desde \textbf{o início do século passado}
      \begin{itemize}
        \item telégrafo sem fio
      \end{itemize}
    \item \textbf{Avanço da tecnologia sem fio}
      \begin{itemize}
        \item rádio e televisão
      \end{itemize}
    \item \textbf{Mais recentemente aparece em}
      \begin{itemize}
        \item telefones celulares
        \item satélites
        \item redes sem fio
      \end{itemize}
  \end{itemize}
  \begin{center}
    \begin{tikzpicture}[font=\scriptsize]
      \node[host](t){telégrafo};
      \node[host,right=1.6cm of t](r){rádio/TV};
      \node[host,right=1.6cm of r](c){celular};
      \node[host,right=1.6cm of c](s){satélite};
      \node[host,right=1.6cm of s](w){Wi-Fi};
      \draw[->,thick,aneliseVerde] (t) -- (r);
      \draw[->,thick,aneliseVerde] (r) -- (c);
      \draw[->,thick,aneliseVerde] (c) -- (s);
      \draw[->,thick,aneliseVerde] (s) -- (w);
      \node[above=2mm of r,font=\tiny]{evolução};
    \end{tikzpicture}\\[1mm]
    {\scriptsize \textbf{Figura 1:} evolução histórica das comunicações sem fio, do telégrafo às redes Wi-Fi.}
  \end{center}
\end{frame}

\begin{frame}{Introdução: aplicações}
  \begin{itemize}
    \item Recebem os \textbf{maiores investimentos em P\&D}
    \item Permitem a \textbf{interconexão de diferentes tipos de dispositivos}
      \begin{itemize}
        \item impressoras, faxes, telefones, agendas eletrônicas, computadores pessoais, IoT\ldots
      \end{itemize}
  \end{itemize}
  \begin{block}{Classificação (segundo a cobertura)}
    \begin{itemize}
      \item \textbf{WPAN} --- redes pessoais sem fio: objetivo eliminar fios (ex.: Bluetooth, ZigBee, UWB)
      \item \textbf{WLAN} --- redes locais sem fio (ex.: IEEE 802.11)
      \item \textbf{WMAN} --- redes metropolitanas sem fio (ex.: IEEE 802.16 WiMAX)
      \item \textbf{WWAN} --- redes de área extensa (celulares 4G/5G, satélite)
    \end{itemize}
  \end{block}
  \begin{itemize}
    \item Várias redes exploram a \textbf{mobilidade do usuário}
      \begin{itemize}
        \item característica importante, intrínseca e exclusiva das redes sem fio
      \end{itemize}
  \end{itemize}
\end{frame}

\begin{frame}{Introdução: meios e problemas}
  \begin{columns}
    \begin{column}{0.48\textwidth}
      \textbf{Transmissão sem fio utiliza o ar como meio}
      \begin{itemize}
        \item radiofrequência
        \item infravermelho
        \item luz (Li-Fi, OWC)
      \end{itemize}
      \vspace{0.2cm}
      \textbf{Principais problemas}
      \begin{itemize}
        \item grande atenuação
        \item limitado alcance
        \item espectro de frequências disponível
        \item interferência
        \item dificuldade de alcançar altas taxas
        \item segurança
        \item reflexões do sinal no interior de uma casa
      \end{itemize}
    \end{column}
    \begin{column}{0.52\textwidth}
      \begin{center}
      \begin{tikzpicture}[font=\scriptsize]
        \node[draw,ellipse,fill=aneliseAzul!12,text width=5.6cm,align=center](esp){Espectro eletromagnético\\ (arco-íris de Maxwell)};
        \node[below=3mm of esp,text width=5.2cm,align=center,font=\scriptsize](uso){Várias faixas podem ser usadas para comunicações};
        \node[draw,fill=aneliseLaranja!15,rounded corners=2mm,below=8mm of uso,text width=6cm,align=center](ism){Bandas ISM (2,4~GHz, 5~GHz, 6~GHz) sem licença\\(desde 1985, FCC: potência $<$ 1~W)};
        \draw[->,dashed,thick] (esp) -- (ism);
      \end{tikzpicture}
      \end{center}
    \end{column}
  \end{columns}
\end{frame}

\section{Conceitos básicos: sinais}

\begin{frame}{Conceitos básicos: sinais}
  \begin{itemize}
    \item Informação transferida da origem ao destino através de \textbf{sinais eletromagnéticos}
    \item Sinais são:
      \begin{itemize}
        \item \textbf{função do tempo} (domínio temporal)
        \item \textbf{função da frequência} (domínio espectral)
      \end{itemize}
  \end{itemize}
  \begin{center}
    \begin{tikzpicture}[font=\scriptsize,scale=0.9]
      \begin{axis}[width=6cm,height=3cm,axis lines=middle,xlabel={$t$},ylabel={$A$},ytick=\empty,xtick=\empty,title={Temporal: sinal periódico}]
        \addplot[domain=0:6.28,samples=100,thick,aneliseAzul]{sin(deg(x)*180/3.14159)};
      \end{axis}
      \begin{axis}[at={(9cm,0cm)},width=6cm,height=3cm,axis lines=middle,xlabel={$f$},ylabel={$A$},ytick=\empty,title={Espectral: picos em $f$, $2f$}]
        \addplot[domain=0:6.28,samples=100,thick,aneliseLaranja]{sin(deg(x)*180/3.14159)};
      \end{axis}
    \end{tikzpicture}
  \end{center}
\end{frame}

\begin{frame}{Domínio do tempo e da frequência}
  \begin{block}{Domínio temporal}
    Um gráfico no domínio do tempo mostra como um sinal varia ao longo do tempo.
    \begin{itemize}
      \item sinais analógicos ou digitais
      \item sinais periódicos ou aperiódicos
    \end{itemize}
  \end{block}
  \begin{block}{Domínio espectral}
    Um sinal eletromagnético pode estar composto de diversas frequências.
    \begin{itemize}
      \item \textbf{Espectro:} faixa de frequências que um sinal contém (ex.: de $f$ a $3f$)
      \item \textbf{Largura de banda absoluta:} largura do espectro (ex.: $3f - f = 2f$)
      \item \textbf{Banda passante efetiva:} banda estreita onde está a maioria da energia
    \end{itemize}
  \end{block}
  \begin{center}
    \scriptsize Descrição temporal $=$ ``como varia com o tempo''; descrição espectral $=$ ``composição em frequências''. (Fonte: Profa. Larissa Driemeier)
  \end{center}
\end{frame}

\section{Regulação espectral}

\begin{frame}{Espectro e regulação: ANATEL}
  \begin{itemize}
    \item Características de um sinal eletromagnético dependem de sua frequência
    \item Diversas faixas podem ser usadas; outras estão reservadas por \textbf{regulação}
  \end{itemize}
  \begin{block}{ANATEL}
    Plano de Atribuição, Distribuição e Destinação de Radiofrequências:
    \begin{itemize}
      \item \url{https://www.anatel.gov.br/setorregulado/atribuicao-destinacao-e-distribucion-de-faixas}
    \end{itemize}
  \end{block}
  \begin{block}{Bandas ISM}
    Desde 1985, o FCC (EUA) permitiu o uso \textbf{sem licença} da banda ISM (Industrial, Scientific, and Medical), sempre que a potência do sinal seja inferior a 1~W. Reguladores equivalentes: ANATEL (BR), ETSI (UE).
  \end{block}
\end{frame}

\begin{frame}{Características do meio de transmissão}
  \begin{columns}
    \begin{column}{0.55\textwidth}
      \begin{itemize}
        \item \textbf{Atenuação} --- potência cai com a distância
        \item \textbf{Perda no espaço livre} --- espalhamento do sinal
        \item \textbf{Ruídos} (térmico, intermodulação, diafonia, impulsivo)
        \item \textbf{Desvanecimento} (\emph{fading}) --- variação temporal da potência
        \item \textbf{Absorção atmosférica} --- vapor de água, oxigênio (picos em 22~GHz e 60~GHz)
      \end{itemize}
      \vspace{0.2cm}
      \textbf{Comparação com o fio:}
      \begin{itemize}
        \item Ethernet: onda \textbf{confinada} $\rightarrow$ menos interferência, menos atenuação; BER $\sim 10^{-7}$
        \item Sem fio: meio \textbf{no-guiado}, BER típica $\sim 10^{-4}$ (até $10^{-2}$ em interiores)
      \end{itemize}
    \end{column}
    \begin{column}{0.45\textwidth}
      \begin{center}
      \begin{tikzpicture}[font=\scriptsize]
        \node[host](Tx){TX};
        \node[host,right=4.5cm of Tx](Rx){RX};
        \draw[thick,dashed,aneliseLaranja] (Tx) -- node[midway,above=2mm,font=\tiny]{sinal direto (atenuado)} (Rx);
        \draw[thick,dashed,aneliseGris] (Tx) to[bend right=35] node[midway,below=1mm,font=\tiny]{refletida} (Rx);
        \draw[thick,dashed,aneliseGris] (Tx) to[bend right=65] node[midway,below=2mm,font=\tiny]{difratada} (Rx);
        \node[below=3mm of Rx,font=\tiny,text width=4cm,align=center]{multipercurso $\rightarrow$ chegada em instantes distintos};
      \end{tikzpicture}
      \end{center}
    \end{column}
  \end{columns}
\end{frame}

\section{Atenuação, ruído, SINR}

\begin{frame}{Atenuação: definição}
  \begin{itemize}
    \item Potência de um sinal cai com a distância $\rightarrow$ \textbf{atenuação}
    \item O sinal recebido deve ter potência suficiente para que o circuito do receptor possa detectá-lo e interpretá-lo
    \item O sinal deve manter um nível suficientemente mais alto que o ruído para ser recebido sem erros
  \end{itemize}
  \begin{block}{Consequências}
    \begin{itemize}
      \item além de uma certa distância, a atenuação é muito forte $\rightarrow$ se usam \textbf{repetidores}
      \item maior no espaço livre (meio no-guiado)
      \item varia com a frequência (maior nas frequências altas) $\rightarrow$ técnicas de \textbf{equalização}
    \end{itemize}
  \end{block}
\end{frame}

\begin{frame}{Perdas no espaço livre e por penetração}
  \begin{columns}
    \begin{column}{0.5\textwidth}
      \textbf{Perda no espaço livre}
      \begin{itemize}
        \item tipo particular de atenuação
        \item o sinal se espalha conforme a distância aumenta $\rightarrow$ atenuação crescente
        \item o sinal se dispersa por uma área cada vez maior
      \end{itemize}
    \end{column}
    \begin{column}{0.5\textwidth}
      \textbf{Perdas por penetração}
      \begin{itemize}
        \item perdas do sinal através de paredes e pisos
        \item normalmente extraídas por medições
      \end{itemize}
      Estudo na rede 802.11 da PUC-Rio:
      \begin{itemize}
        \item espaço livre: 0~dB
        \item janelas (tinta metálica): 3--8~dB
        \item paredes finas 5--8~dB; médias 10~dB; espessas (15~cm) 15--20~dB; muito espessas (30~cm) 20--25~dB
        \item pisos e tetos: 15--25~dB
      \end{itemize}
    \end{column}
  \end{columns}
\end{frame}

\begin{frame}{Ruídos}
  \textbf{Alterações sofridas pelo sinal transmitido entre transmissão e recepção.} Quatro tipos principais:
  \begin{description}
    \item[Térmico] Uniformemente distribuído através do espectro
    \item[Intermodulação] Compartilhamento de um meio entre sinais de diferentes frequências
    \item[Diafonia (\emph{crosstalk})] Sinais indesejados são ``induzidos'' no sinal recebido
    \item[Impulsivo] Pulsos irregulares de curta duração com relativamente grande amplitude
  \end{description}
\end{frame}

\begin{frame}{Razão Sinal-Ruído (SINR)}
  \begin{block}{Definição}
    Usada para descrever a qualidade do sinal: medida do nível do sinal relativo ao ruído de fundo.
    Quanto maior o SINR, melhor a qualidade.
    \begin{itemize}
      \item normalmente fornecido em \textbf{decibel (dB)}
      \item $\mathrm{SINR} = 10\log_{10}(S/R)$, com $S=$ potência do sinal, $R=$ potência do ruído
    \end{itemize}
  \end{block}
  \begin{block}{Condições para correta recepção}
    \begin{enumerate}
      \item Nível do sinal acima do \textbf{limiar de recepção}: $S >$ threshold (W)
      \item \textbf{SINR acima do limiar de captura}: $10\log_{10}(S/R) >$ capture threshold (dB)
    \end{enumerate}
  \end{block}
  \note{Armadilha: SINR é razão de potências em escala logarítmica; lembre que +10~dB corresponde a 10$\times$ em potência.}
\end{frame}

\section{Características do enlace sem fio}

\begin{frame}{Diferenças em relação ao enlace cabeado}
  \begin{enumerate}
    \item \textbf{Força reduzida do sinal:} os sinais de rádio se atenuam ao se propagar através da matéria (\emph{path loss})
    \item \textbf{Interferência de outras fontes:} as frequências padronizadas (ex.: 2,4~GHz) são compartilhadas por outros equipamentos (telefones sem fio, microondas, motores)
    \item \textbf{Propagação de múltiplos caminhos:} o sinal se reflete no solo e em objetos; o sinal principal e os refletidos chegam em instantes ligeiramente diferentes
  \end{enumerate}
  \begin{alertblock}{Conclusão}
    Tudo isso faz a comunicação através de enlaces sem fio (até ponto-a-ponto) muito mais ``difícil''.
  \end{alertblock}
\end{frame}

\begin{frame}{Múltiplos caminhos: mecanismos de propagação}
  \begin{columns}
    \begin{column}{0.32\textwidth}
      \textbf{Reflexão}
      \begin{itemize}
        \item superfície grande comparada ao comprimento de onda
        \item sinais podem ser recebidos sem visada direta
      \end{itemize}
    \end{column}
    \begin{column}{0.32\textwidth}
      \textbf{Difração}
      \begin{itemize}
        \item canto de corpo impenetrável \textbf{grande} comparado ao comprimento de onda
        \item se propaga em diferentes direções
        \item sem visada direta
      \end{itemize}
    \end{column}
    \begin{column}{0.32\textwidth}
      \textbf{Dispersão (espalhamento)}
      \begin{itemize}
        \item canto/corpo \textbf{pequeno} ou da mesma ordem do comprimento de onda
        \item se propaga em várias direções
      \end{itemize}
    \end{column}
  \end{columns}
  \vspace{0.3cm}
  \begin{center}
    \begin{tikzpicture}[font=\scriptsize]
      \node[host](Tx){TX};
      \node[host,right=6cm of Tx](Rx){RX};
      \draw[->,thick,aneliseAzul] (Tx) -- (Rx);
      \node[draw,fill=aneliseGris!30,rounded corners=1mm,below=1.2cm of Tx](obs){obstáculo};
      \draw[->,dashed,thick,aneliseLaranja] (Tx.south) -- (obs.north);
      \draw[->,dashed,thick,aneliseLaranja] (obs.east) to[bend right=25] (Rx.south);
      \node[below=2mm of obs,font=\tiny]{reflexão/difração/dispersão};
    \end{tikzpicture}
  \end{center}
\end{frame}

\begin{frame}{Desvanecimento (\emph{fading})}
  \begin{block}{Definição}
    Variação temporal da potência do sinal recebido, causada por mudanças no meio ou no(s) caminho(s).
  \end{block}
  \begin{columns}
    \begin{column}{0.5\textwidth}
      \textbf{Ambiente fixo}
      \begin{itemize}
        \item afetado por mudanças atmosféricas (ex.: chuva)
      \end{itemize}
      \textbf{Ambiente móvel}
      \begin{itemize}
        \item afetado pela localização relativa de obstáculos
      \end{itemize}
    \end{column}
    \begin{column}{0.5\textwidth}
      Causado por reflexão, difração, espalhamento:
      \begin{itemize}
        \item reflexão: interferência construtiva/destrutiva no receptor
        \item difração: obstrução por superfícies irregulares
        \item espalhamento: muitos objetos pequenos
      \end{itemize}
      Flutuações: \emph{shadowing}, pequenas mudanças de posição/endereço da antena, obstáculos móveis.
    \end{column}
  \end{columns}
\end{frame}

\begin{frame}{Absorção atmosférica}
  \begin{columns}
    \begin{column}{0.55\textwidth}
      \textbf{Vapor de água}
      \begin{itemize}
        \item pico de atenuação em 22~GHz
        \item menor atenuação abaixo de 15~GHz
      \end{itemize}
      \textbf{Oxigênio}
      \begin{itemize}
        \item pico em 60~GHz
        \item menor atenuação abaixo de 30~GHz
      \end{itemize}
      \textbf{Chuva e névoa}
      \begin{itemize}
        \item dispersam ondas $\rightarrow$ atenuação (causa principal de perdas em zonas chuvosas)
        \item soluções: distâncias curtas ou bandas de baixa frequência
      \end{itemize}
    \end{column}
    \begin{column}{0.45\textwidth}
      \begin{center}
      \begin{tikzpicture}[font=\scriptsize,
          camada/.style={draw,fill=aneliseAzul!10,text width=4.4cm,align=center,rounded corners=1mm,inner sep=1mm},
        ]
        \node[camada](v){Vapor d'água\\ pico 22~GHz};
        \node[camada,below=4mm of v](o){Oxigênio\\ pico 60~GHz};
        \node[camada,below=4mm of o](c){Chuva/névoa\\ dispersão};
      \end{tikzpicture}
      \end{center}
    \end{column}
  \end{columns}
\end{frame}

\begin{frame}{Modelos de propagação}
  Um sinal irradiado pode se propagar de três formas:
  \begin{center}
    \begin{tikzpicture}[font=\scriptsize]
      \node[host](Tx){TX};
      \node[host,right=6cm of Tx](Rx){RX};
      % ground wave
      \draw[thick,aneliseVerde] (Tx) to[bend right=60] node[midway,below=3mm,font=\tiny]{acima do solo ($\leq$3~MHz)} (Rx);
      % sky wave
      \draw[thick,aneliseLaranja] (Tx) to[bend left=90] node[midway,above=3mm,font=\tiny]{no céu (3--30~MHz, ionosfera)} (Rx);
      % line of sight
      \draw[thick,aneliseAzul] (Tx) -- node[midway,above=1mm,font=\tiny]{visada direta ($>$30~MHz)} (Rx);
    \end{tikzpicture}
  \end{center}
  \begin{itemize}
    \item \textbf{Acima do solo} (\emph{ground wave}): segue o contorno da terra; até 3~MHz; ex.: rádio AM
    \item \textbf{No céu} (\emph{sky wave}): refração na ionosfera, saltos terra-ionosfera; 3--30~MHz; ex.: rádio amador
    \item \textbf{Visada direta} (\emph{line-of-sight}): antenas alinhadas; satélite e comunicações terrestres (refração)
  \end{itemize}
\end{frame}

\begin{frame}{Efeitos da visada direta}
  \begin{itemize}
    \item \textbf{Com visada direta:}
      \begin{itemize}
        \item difração e dispersão têm efeitos menores
        \item reflexão pode ter impacto significativo
      \end{itemize}
    \item \textbf{Sem visada direta:}
      \begin{itemize}
        \item difração e dispersão são os meios primários de recepção
      \end{itemize}
  \end{itemize}
\end{frame}

\section{Multiplexação e duplexação}

\begin{frame}{Multiplexação e duplexação}
  \begin{columns}
    \begin{column}{0.5\textwidth}
      \textbf{Multiplexação} (compartilhar o meio físico)
      \begin{itemize}
        \item divisão do meio na camada física
        \item centralizada em um \textbf{multiplexador}
        \item classificação segundo a variável de separação:
          \begin{itemize}
            \item \textbf{TDM} (tempo) --- usada em IEEE 802.16
            \item \textbf{FDM} (frequência)
            \item \textbf{WDM} (comprimento de onda, fibra)
          \end{itemize}
      \end{itemize}
    \end{column}
    \begin{column}{0.5\textwidth}
      \textbf{Duplexação} (comunicação entre duas estações)
      \begin{itemize}
        \item \textbf{simplex:} um único sentido
        \item \textbf{half-duplex}: dois sentidos, mas não simultâneos
        \item \textbf{full-duplex:} dois sentidos simultâneos
      \end{itemize}
      \begin{center}
        \begin{tikzpicture}[font=\scriptsize]
          \node[host](A){A};
          \node[host,right=2.6cm of A](B){B};
          \draw[->,thick,aneliseAzul] (A) -- node[midway,above=1mm,font=\tiny]{simplex} (B);
          \draw[<->,thick,aneliseLaranja] ([yshift=1mm]A.south) -- node[midway,below=1mm,font=\tiny]{half-duplex} ([yshift=1mm]B.south);
          \draw[<->,thick,aneliseVerde] ([yshift=3mm]A.south) to[bend left=20] node[midway,below=4mm,font=\tiny]{full-duplex} ([yshift=3mm]B.south);
        \end{tikzpicture}
      \end{center}
    \end{column}
  \end{columns}
\end{frame}

\section{Técnicas de transmissão}

\begin{frame}{Codificação e modulação}
  \begin{block}{Codificações (\emph{banda base})}
    \begin{itemize}
      \item trabalham em banda base
      \item preservam as faixas de frequências originais dos dados
      \item ex.: PPM (pulse position modulation), NRZ, Manchester
    \end{itemize}
  \end{block}
  \begin{block}{Modulações (\emph{banda passante})}
    \begin{itemize}
      \item modificam a faixa de frequências dos dados
      \item usam uma portadora de alta frequência (senoide)
      \item variação de amplitude, frequência ou fase $\rightarrow$ ASK, FSK, PSK, QAM
    \end{itemize}
  \end{block}
  \begin{alertblock}{Por que modular?}
    Ondas quadradas das codificações contêm múltiplas frequências; atenuação e velocidade de propagação variam com a frequência $\rightarrow$ codificação em banda base só para curtas distâncias e baixas velocidades. Alguns meios (fibra, RF) exigem que o sinal ocupe uma determinada faixa; além disso, múltiplas estações devem acessar o mesmo meio físico.
  \end{alertblock}
\end{frame}

\begin{frame}{Modulações: ASK, FSK, PSK}
  \begin{columns}
    \begin{column}{0.33\textwidth}
      \textbf{ASK --- Amplitude Shift Keying}
      \begin{itemize}
        \item M-ASK: varia a amplitude entre $M$ valores
        \item BASK (M=2): duas amplitudes
      \end{itemize}
    \end{column}
    \begin{column}{0.33\textwidth}
      \textbf{FSK --- Frequency Shift Keying}
      \begin{itemize}
        \item M-FSK: varia a frequência entre $M$ valores
        \item BFSK (M=2): duas frequências
      \end{itemize}
    \end{column}
    \begin{column}{0.33\textwidth}
      \textbf{PSK --- Phase Shift Keying}
      \begin{itemize}
        \item M-PSK: varia a fase entre $M$ valores
        \item BPSK: transição $\rightarrow$ fase invertida $180^\circ$
      \end{itemize}
    \end{column}
  \end{columns}
  \vspace{0.3cm}
  \begin{center}
    \begin{tikzpicture}[font=\scriptsize,scale=0.9]
      \begin{axis}[width=5cm,height=2.2cm,axis lines=middle,ytick=\empty,xtick=\empty,xlabel={$t$},title={ASK}]
        \addplot[domain=0:6.28,samples=120,thick,aneliseAzul]{1.2*sin(deg(x)*180/3.14159)*(sin(deg(2*x)*180/3.14159)>0?1:0.2)};
      \end{axis}
      \begin{axis}[at={(6.5cm,0cm)},width=5cm,height=2.2cm,axis lines=middle,ytick=\empty,xtick=\empty,xlabel={$t$},title={FSK}]
        \addplot[domain=0:6.28,samples=120,thick,aneliseLaranja]{sin(deg(x)*(2+0.8*(sin(deg(x)*180/3.14159)>0?1:0))*180/3.14159)};
      \end{axis}
      \begin{axis}[at={(13cm,0cm)},width=5cm,height=2.2cm,axis lines=middle,ytick=\empty,xtick=\empty,xlabel={$t$},title={BPSK (transição de fase)}]
        \addplot[domain=0:6.28,samples=120,thick,aneliseVerde]{sin(deg(x)*180/3.14159 + (x>3.14?3.14159:0))};
      \end{axis}
    \end{tikzpicture}
  \end{center}
\end{frame}

\begin{frame}{QAM --- Quadrature Amplitude Modulation}
  \begin{columns}
    \begin{column}{0.52\textwidth}
      \begin{itemize}
        \item M-QAM: varia \textbf{amplitude e fase} da portadora entre $M$ valores
        \item quanto maior o número de símbolos, maior a probabilidade de erros (ruído, atenuação, dificuldade de distinguir símbolos)
        \item usada no IEEE 802.11 e 802.16; base de Wi-Fi 6 (1024-QAM, 4096-QAM)
      \end{itemize}
      \begin{exampleblock}{Exemplo (OFDM 802.11a/g)}
        $6$ bits/símbolo $\times 48$ subcanais $= 288$ bits/channel; $250\,000$ símbolos/s $\rightarrow$ $250\,000 \times 216 = 54$~Mbps
      \end{exampleblock}
    \end{column}
    \begin{column}{0.48\textwidth}
      \begin{center}
        \begin{tikzpicture}[font=\scriptsize,scale=0.85]
          \begin{axis}[width=5.6cm,height=5.4cm,axis lines=middle,xlabel={I},ylabel={Q},title={Constelação 16-QAM},xtick=\empty,ytick=\empty,ymin=-2.2,ymax=2.2,xmin=-2.2,xmax=2.2]
            \addplot[only marks,mark=*,mark size=3pt,thick,aneliseAzul] coordinates {
              (-1.5,-1.5)(-0.5,-1.5)(0.5,-1.5)(1.5,-1.5)
              (-1.5,-0.5)(-0.5,-0.5)(0.5,-0.5)(1.5,-0.5)
              (-1.5,0.5)(-0.5,0.5)(0.5,0.5)(1.5,0.5)
              (-1.5,1.5)(-0.5,1.5)(0.5,1.5)(1.5,1.5)
            };
          \end{axis}
        \end{tikzpicture}
        \\ \scriptsize Constelação 16-QAM (4 bits/símbolo)
      \end{center}
    \end{column}
  \end{columns}
  \note{Relacione a constelação com a taxa de erro: mais bits por símbolo exigem melhor SINR --- este é o elo com o orçamento de enlace da aula.}
\end{frame}

\section{Espalhamento de espectro}

\begin{frame}{Spread Spectrum: motivação}
  \begin{block}{Ideia}
    Espalhar um sinal em uma largura de banda superior à exigida pelo sinal original.
  \end{block}
  \textbf{Objetivos:}
  \begin{itemize}
    \item minimizar interferência em frequências específicas
    \item minimizar interferência por múltiplos caminhos
  \end{itemize}
  \textbf{Custos:} desperdício de banda para um usuário individual.

  \textbf{Vantagem:} vários usuários (redes) podem usar simultaneamente a mesma banda sem interferências significativas $\rightarrow$ eficiente em termos de banda.
  \begin{block}{Sequências pseudo-aleatórias}
    Parecem aleatórias, mas podem ser reproduzidas pelos receptores (chip sequence).
  \end{block}
  \begin{itemize}
    \item Interferências são (na sua maioria) de \textbf{banda estreita} $\rightarrow$ interferem só numa pequena faixa do sinal
    \item Minimiza atenuação por múltiplos caminhos; permite compartilhar faixa; privacidade
    \item 1999: OFDM e HR-DSSS para maior largura de banda
  \end{itemize}
\end{frame}

\begin{frame}{FHSS --- Frequency Hopping Spread Spectrum}
  \begin{columns}
    \begin{column}{0.55\textwidth}
      \begin{itemize}
        \item modulação com portadora que \textbf{salta de frequência} em função do tempo sobre uma larga faixa
        \item o \textbf{código de salto} (hopping code) determina as frequências e a ordem
        \item recomendação FCC banda 2,4~GHz: $\geq 75$ frequências por canal, máximo 400~ms por salto
        \item um sinal de banda estreita só interfere se transmite na mesma frequência ao mesmo tempo
        \item máximo 2--3~Mbps (acima disso, aumentam os erros)
      \end{itemize}
    \end{column}
    \begin{column}{0.45\textwidth}
      \begin{exampleblock}{Características}
        \begin{itemize}
          \item permite coexistência de várias redes na mesma área (sequências de saltos distintas)
          \item usado em Bluetooth e primeiras versões 802.11
        \end{itemize}
      \end{exampleblock}
      \begin{center}
        \begin{tikzpicture}[font=\scriptsize]
          \node[host](a){BT};
          \node[host,right=1.6cm of a](b){BT};
          \node[host,right=1.6cm of b](c){Wi-Fi};
          \draw[thick,aneliseLaranja] (a) to[bend left=20] (b);
          \draw[thick,aneliseVerde] (b) -- (c);
        \end{tikzpicture}
      \end{center}
    \end{column}
  \end{columns}
\end{frame}

\begin{frame}{DSSS --- Direct Sequence Spread Spectrum}
  \begin{itemize}
    \item multiplexação do sinal por uma \textbf{sequência pseudo-aleatória} (chipping) de taxa muito maior
    \item XOR dos dados com a sequência: bit $0 \rightarrow$ sequência; bit $1 \rightarrow$ complemento
    \item o \textbf{fator de espalhamento} (número de chips por bit) determina a banda passante do sinal resultante
    \item 802.11b: sequência de 11 chips $\rightarrow$ 11~Mbps
  \end{itemize}
  \begin{block}{Trade-off}
    \begin{itemize}
      \item mais caro e maior consumo de energia
      \item maior robustez frente a interferência de banda estreita e multipercurso
    \end{itemize}
  \end{block}
\end{frame}

\begin{frame}{OFDM --- Orthogonal Frequency Division Multiplexing}
  \begin{itemize}
    \item divide o espectro em \textbf{múltiplas sub-bandas estreitas}; envia bits em cada sub-banda
    \item todas as sub-bandas dedicadas a uma única fonte (diferente do FDM)
    \item subportadoras \textbf{ortogonais} (centradas nos zeros das adjacentes)
  \end{itemize}
  \begin{block}{802.11a/g/n}
    \begin{itemize}
      \item canal de 20~MHz, 52 subportadoras (48 dados + 4 controle)
      \item transmissão em paralelo; canais sobrepostos sem interferência mútua
      \item taxa de 250\,000 símbolos/s
    \end{itemize}
  \end{block}
  \begin{exampleblock}{Evolução (atualizado)}
    802.11n: OFDM 108 subportadoras, MIMO 4$\times$4\\802.11ac: OFDM 234, QAM-256, MIMO 8$\times$8\\802.11ax (Wi-Fi 6): OFDMA, QAM-1024, até 9,6~Gbps
  \end{exampleblock}
\end{frame}

\begin{frame}{Padrões 802.11 (resumo)}
  \begin{center}
    \scriptsize
    \begin{tabular}{@{}llll@{}}
      \toprule
      \textbf{Padrão} & \textbf{Banda (MHz)} & \textbf{Técnica} & \textbf{Observações}\\
      \midrule
      802.11 & FHSS/DSSS, BPSK & 2,4 & 1--2 Mbps\\
      802.11b & DSSS, BPSK/QPSK & 2,4 & 11 Mbps (chip de 11)\\
      802.11a/g & OFDM (52 subport.), QAM-64 & 5 / 2,4 & 54 Mbps\\
      802.11n & OFDM, QAM-64, MIMO 4$\times$4 & 2,4/5 & 600 Mbps\\
      802.11ac & OFDM 234, QAM-256, MIMO 8$\times$8 & 5 & até 6,9 Gbps\\
      802.11ax (Wi-Fi 6) & OFDMA, QAM-1024, MIMO multi-user & 2,4/5/6 & até 9,6 Gbps\\
      \bottomrule
    \end{tabular}
  \end{center}
  \vspace{0.2cm}
  \begin{alertblock}{Nota atualizada}
    Wi-Fi 6E (802.11ax na banda de 6~GHz) e Wi-Fi 7 (802.11be, 2024, 320~MHz, QAM-4096, até ~46~Gbps) são os padrões atuais.
  \end{alertblock}
\end{frame}

\begin{frame}{Vantagens de OFDM}
  \begin{itemize}
    \item \textbf{Maior imunidade a ruído e interferências}
      \begin{itemize}
        \item subportadoras independentes: pode-se descartar uma ou várias e usar as restantes
      \end{itemize}
    \item usado em IEEE 802.11 e 802.16 (WiMAX)
  \end{itemize}
\end{frame}

\section{Redes com e sem infraestrutura}

\begin{frame}{Redes sem fio: duas categorias}
  \begin{columns}
    \begin{column}{0.5\textwidth}
      \textbf{Redes com infraestrutura}
      \begin{itemize}
        \item toda a comunicação através de um \textbf{ponto de acesso}
        \item ex.: redes celulares, Wi-Fi em modo infraestrutura
      \end{itemize}
      \begin{center}
        \begin{tikzpicture}[font=\scriptsize]
          \node[host,draw,circle,fill=aneliseLaranja!25](ap){AP};
          \node[host,left=1.4cm of ap](a){};
          \node[host,right=1.4cm of ap](b){};
          \node[host,below=1.2cm of ap](c){};
          \draw[thick] (a) -- (ap);
          \draw[thick] (b) -- (ap);
          \draw[thick] (c) -- (ap);
        \end{tikzpicture}
      \end{center}
    \end{column}
    \begin{column}{0.5\textwidth}
      \textbf{Redes sem infraestrutura (ad hoc)}
      \begin{itemize}
        \item estações se comunicam diretamente
        \item ad hoc de comunicação direta
        \item ad hoc de múltiplos saltos (nós também roteadores)
      \end{itemize}
      \begin{center}
        \begin{tikzpicture}[font=\scriptsize]
          \node[host](a){};
          \node[host,right=1.6cm of a](b){};
          \node[host,right=1.6cm of b](c){};
          \draw[thick] (a) -- (b) node[midway,above=1mm,font=\tiny]{direta};
          \draw[thick] (b) -- (c) node[midway,above=1mm,font=\tiny]{multi-salto};
        \end{tikzpicture}
      \end{center}
    \end{column}
  \end{columns}
\end{frame}

\begin{frame}{Redes ad hoc: características e aplicações}
  \begin{columns}
    \begin{column}{0.55\textwidth}
      \textbf{Características}
      \begin{itemize}
        \item auto-organização dinâmica
        \item topologia arbitrária e temporal
      \end{itemize}
      \textbf{Vantagens}
      \begin{itemize}
        \item grande flexibilidade
        \item baixo custo de instalação
        \item robustez
      \end{itemize}
      \textbf{Aplicações}
      \begin{itemize}
        \item ambientes sem infraestrutura ou com infraestrutura não confiável
        \item reuniões inesperadas, treinos militares, catástrofes, drones (FANET)
      \end{itemize}
    \end{column}
    \begin{column}{0.45\textwidth}
      \begin{block}{Por que não CSMA/CD?}
        \begin{itemize}
          \item grande diferença de potência entre sinal transmitido e recebido
          \item difícil separação sinal/ruído, transmissão/recepção
          \item nem todas as estações escutam as outras (colisão no transmissor não implica colisão no receptor)
          \item problema do \textbf{terminal escondido}
        \end{itemize}
      \end{block}
    \end{column}
  \end{columns}
\end{frame}

\begin{frame}{Terminal escondido e terminal exposto}
  \begin{columns}
    \begin{column}{0.5\textwidth}
      \textbf{Terminal escondido}
      \begin{center}
        \begin{tikzpicture}[font=\scriptsize]
          \node[host](A){A};
          \node[host,right=1.6cm of A](B){B};
          \node[host,right=1.6cm of B](C){C};
          \draw[->,thick,aneliseAzul] (A) -- node[midway,above=1mm,font=\tiny]{TX} (B);
          \draw[->,thick,aneliseLaranja] (C) -- node[midway,below=1mm,font=\tiny]{TX} (B);
          \node[below=2mm of B,font=\tiny,text width=3cm,align=center]{A e C não se ouvem $\rightarrow$ colisão em B};
        \end{tikzpicture}
      \end{center}
      \begin{itemize}
        \item ocorre com sobreposição de áreas de cobertura
        \item não importa a interferência no transmissor, e sim no receptor
        \item protocolos centralizados ou coordenação de APs ou canais distintos
      \end{itemize}
    \end{column}
    \begin{column}{0.5\textwidth}
      \textbf{Terminal exposto}
      \begin{center}
        \begin{tikzpicture}[font=\scriptsize]
          \node[host](A){A};
          \node[host,right=1.4cm of A](B){B};
          \node[host,right=1.4cm of B](C){C};
          \node[host,right=1.4cm of C](D){D};
          \draw[->,thick,aneliseAzul] (A) -- (B);
          \draw[->,thick,aneliseAzul] (C) -- (D);
          \node[below=2mm of C,font=\tiny,text width=3.4cm,align=center]{C pode transmitir a D, mas ``pensa que não''};
        \end{tikzpicture}
      \end{center}
      \begin{itemize}
        \item C escuta a A $\rightarrow$ adia; mas D está fora da área de A
        \item desperdício de oportunidade (CSMA puro não funciona bem)
      \end{itemize}
    \end{column}
  \end{columns}
\end{frame}

% ---------------------------------------------------------------------------
\begin{frame}{Exercício resolvido: orçamento de enlace}
  \textbf{Problema:} Um ponto de acesso transmite com potência $P_t = 20$~dBm e ganho de antena $G_t = 5$~dBi; o receptor tem ganho $G_r = 5$~dBi. Os equipamentos estão separados por $d = 1$~km, operando em $f = 2{,}4$~GHz. Calcule a perda de percurso no espaço livre e a potência recebida. (Sensibilidade do receptor: $-85$~dBm.)
  \pause
  \begin{block}{Solução passo a passo}
    \begin{align*}
      PL &= 20\log_{10}(d) + 20\log_{10}(f) + 32{,}44
      = 20\log_{10}(1) + 20\log_{10}(2400) + 32{,}44
      \approx 100~\text{dB},\\[2mm]
      P_r &= P_t + G_t + G_r - PL
      = 20 + 5 + 5 - 100 = \alert{-70~\text{dBm}}.
    \end{align*}
    \textbf{Verificação:} $-70 > -85$~dBm $\Rightarrow$ sinal recebido fica \alert{15~dB acima da sensibilidade} (margem de 15~dB).
  \end{block}
  \pause
  \textbf{Conclusão:} a recepção é viável; a margem de 15~dB pode absorver atenuações extras (chuva, obstáculos).
\end{frame}

% ---------------------------------------------------------------------------
\begin{frame}{Resumo da aula}
  \begin{itemize}
    \item Redes sem fio oferecem \textbf{mobilidade, flexibilidade e baixo custo de instalação}, mas têm \textbf{menor QoS}, interferência e restrições regulatórias.
    \item Sinais são descritos no \textbf{domínio do tempo} e da \textbf{frequência}; o espectro é regulado (ANATEL, FCC, ETSI) e as bandas ISM dispensam licença.
    \item O meio sem fio sofre \textbf{atenuação, ruído, desvanecimento e propagação por múltiplos caminhos} (reflexão, difração, dispersão).
    \item O \textbf{SINR} e o \textbf{orçamento de enlace} determinam se a recepção é correta.
    \item \textbf{Multiplexação} (TDM/FDM/WDM) e \textbf{duplexação} (simplex, half/full-duplex) organizam o compartilhamento do meio.
    \item \textbf{FHSS, DSSS e OFDM} são as principais técnicas de espalhamento espectral; OFDM é a base de 802.11a/g/n/ac/ax.
    \item Redes podem ser \textbf{com infraestrutura} (AP) ou \textbf{ad hoc}; o acesso ao meio sofre com os problemas de \textbf{terminal escondido e exposto}.
  \end{itemize}
\end{frame}

% ---------------------------------------------------------------------------
\begin{frame}{Autoavaliação}
  \begin{enumerate}
    \item O SINR é normalmente expresso em:
      \begin{itemize}
        \item (a) watts
        \item (b) decibéis (dB)
        \item (c) bits por segundo
      \end{itemize}
      \pause
      \textbf{Gabarito:} (b)

    \item Qual técnica de espalhamento espectral é a base do 802.11a/g/n/ac/ax?
      \begin{itemize}
        \item (a) FHSS
        \item (b) DSSS
        \item (c) OFDM
      \end{itemize}
      \pause
      \textbf{Gabarito:} (c)

    \item Se $P_t = 20$~dBm, $G_t = G_r = 5$~dBi e $PL = 100$~dB, a potência recebida é:
      \begin{itemize}
        \item (a) $-70$~dBm
        \item (b) $+130$~dBm
        \item (c) $-100$~dBm
      \end{itemize}
      \pause
      \textbf{Gabarito:} (a)
  \end{enumerate}
\end{frame}

% ---------------------------------------------------------------------------
\begin{frame}{Laboratório sugerido: associação 802.11 com Wireshark}
  \textbf{Objetivo de aprendizagem:} observar o handshake de associação e os quadros de gerenciamento 802.11 em modo monitor.
  \begin{enumerate}
    \item Verifique se a interface sem fio suporta \textbf{modo monitor} (Linux/macOS: \texttt{iw} ou \texttt{airmon-ng}; Windows: adaptador compatível com npcap em modo monitor).
    \item Desative a associação automática (\emph{auto-join}) e coloque a interface em modo monitor na banda 2,4~GHz.
    \item Inicie a captura no Wireshark e \textbf{associe-se} à rede (ou apenas observe as redes vizinhas).
    \item Filtre por \texttt{wlan.fc.type == 0} (quadros de gerenciamento) e identifique \texttt{Probe Request/Response}, \texttt{Authentication} e \texttt{Association Request/Response}.
    \item Observe o \emph{handshake} 4-way do WPA2 (filtro \texttt{eapol}) e as capacidades HT/VHT negociadas.
    \item Compare o comportamento observado com os problemas de terminal escondido/exposto discutidos em aula.
  \end{enumerate}
\end{frame}

\section{Bibliografia}

\begin{frame}{Bibliografia}
  \begin{itemize}
    \item \textbf{\kurose{}} --- Capítulo 6 (Redes sem fio e redes móveis)
    \item \textbf{\forouzan{}} --- Capítulos 3, 4, 5
    \item \textbf{\stallings{}} --- Capítulos 1, 2, 5, 6, 7 e 11
    \item \textbf{\tanenbaum{}} --- Capítulos 1, 2 e 4
  \end{itemize}
  \begin{block}{Bibliografia adicional}
    \begin{itemize}
      \item \emph{An Introduction to Direct-Sequence Spread-Spectrum Communications} (Maxim AN1890) --- \url{https://www.analog.com/en/resources/technical-articles/introduction-to-spreadspectrum-communications--maxim-integrated.html}
      \item Andersen, Rappaport, Yoshida, ``Propagation Measurements and Models for Wireless Communications Channels'', \emph{IEEE Communications Magazine}, Jan.~1995
      \item Jurdak, Lopes, Baldi, ``A Survey Classification and Comparative Analysis of MAC Protocols for Ad Hoc Networks'', \emph{IEEE Comm.\ Surveys \& Tutorials}, 2004
      \item Xu, Saadawi, ``Does the IEEE 802.11 MAC protocol work well in multihop wireless adhoc networks?'', \emph{IEEE Communications Magazine}, Jun.~2001
    \end{itemize}
  \end{block}
\end{frame}

% ---------------------------------------------------------------------------
\begin{frame}{Leitura recomendada}
  \begin{itemize}
    \item \textbf{\kurose{}} --- cap.~6 (redes sem fio e redes móveis)
    \item \textbf{\forouzan{}} --- caps.~3--5 (sinais, transmissão e meios)
    \item \textbf{\stallings{}} --- caps.~1, 2, 5--7 (fundamentos e propagação)
    \item \textbf{\tanenbaum{}} --- caps.~1, 2 e 4
  \end{itemize}
  \begin{exampleblock}{Dica de estudo}
    Associe cada técnica de espalhamento (FHSS, DSSS, OFDM) aos padrões 802.11 correspondentes da tabela-resumo.
  \end{exampleblock}
\end{frame}

\end{document}